\problem{6.35}{A Galilean boost as a unitary quantum transformation}
Define
\[
G(v,t)=\exp\left[
\frac{\ii}{\hb}(mv\x-vt\p)\right]
\]
with the textbook's active-state
boost convention.
\textbf{(a)} For infinitesimal $v=\delta$,
write $G=I+A+O(\delta^2)$,
$A=\ii\delta(mx-tp)/\hb$.
As $G^\dagger xG=x+[x,A]+O(\delta^2)$,
the canonical commutator gives
\[
[x,A]=\delta t,\qquad
[p,A]=m\delta.
\]
Therefore
\begin{equation}
\boxed{G^\dagger xG=x+\delta t,\qquad
G^\dagger pG=p+m\delta+O(\delta^2).}
\end{equation}
The wave packet's position and
momentum are shifted according to
a frame moving with relative speed $v$.
\textbf{(b)} Let
$A=\ii mvx/\hb$, $B=-\ii vt p/\hb$.
Their commutator is the scalar
$[A,B]=\ii mv^2t/\hb$, and all
nested commutators vanish.
The Baker--Campbell--Hausdorff
identity therefore yields
\begin{align}
G&=\ee^{A+B}
=\ee^{A}\ee^{B}\ee^{-[A,B]/2}\\
&=\boxed{
\exp\left[\frac{\ii}{\hb}
(mvx-\tfrac12mv^2t)\right]\T(vt)},\\
G&=\ee^{B}\ee^{A}\ee^{+[A,B]/2}\\
&=\boxed{\T(vt)
\exp\left[\frac{\ii}{\hb}
(mvx+\tfrac12mv^2t)\right]}.
\end{align}
\textbf{(c)} If $\Psi$ satisfies
$\ii\hb\partial_t\Psi=
[-\hb^2\partial_x^2/(2m)+V(x)]\Psi$,
then the boosted wave is
\begin{equation}
\boxed{\Psi'(x,t)=
\ee^{\ii(mvx-mv^2t/2)/\hb}\,
\Psi(x-vt,t).}
\end{equation}
Set $y=x-vt$ and denote the phase
by $F$. The time derivative is
\[
\ii\hb\partial_t\Psi'
=F[\tfrac12mv^2\Psi-\ii\hb v\partial_y\Psi+
\ii\hb\partial_t\Psi].
\]
Differentiate twice in $x$:
\[
-\frac{\hb^2}{2m}\partial_x^2\Psi'
=F[\tfrac12mv^2\Psi-\ii\hb v\partial_y\Psi
-\tfrac{\hb^2}{2m}\partial_y^2\Psi].
\]
Insert the original Schrödinger
equation; terms match term by term,
and the remaining potential is
$V(y)=V(x-vt)$:
\begin{equation}
\boxed{\ii\hb\partial_t\Psi'
=\left[-\frac{\hb^2}{2m}\partial_x^2
+V(x-vt)\right]\Psi'.}
\end{equation}
\textbf{(d)} This reproduces the
transformation of a wave packet
under a uniform Galilean boost
in Problem 2.50(a), including its
necessary space--time dependent
phase. Simply replacing
$x\mapsto x-vt$ without that phase
would not solve the moving-potential
Schrödinger equation.

\problem{6.36}{Spinless time reversal and the energy spectrum}
For a spinless particle in
position representation,
$\Theta=K$ (complex conjugation),
which is \emph{antiunitary}.
\textbf{(a)} Apply it to an arbitrary
test function $f$:
\[
(\Theta^{-1}x\Theta f)(x)=
[x f^*(x)]^*=xf(x).
\]
For momentum,
\[
(\Theta^{-1}p\Theta f)(x)
=[-\ii\hb (f^*)'(x)]^*
=+\ii\hb f'(x)=-pf(x).
\]
Thus
\begin{equation}
\boxed{\Theta^{-1}x\Theta=x,\qquad
\Theta^{-1}p\Theta=-p.}
\end{equation}
\textbf{(b)} Reversing motion should
undo unitary evolution:
\[
\Theta U(t)\Theta^{-1}=U(-t).
\]
Expand the operator
$U(\delta)=I-\ii H\delta/\hb+O(\delta^2)$.
Because $\Theta$ is antiunitary,
$\Theta\ii\Theta^{-1}=-\ii$.
Hence
\[
\Theta U(\delta)\Theta^{-1}
=I+\frac{\ii\delta}{\hb}
\Theta H\Theta^{-1}+O(\delta^2).
\]
Compare with
$U(-\delta)=I+\ii H\delta/\hb+O(\delta^2)$,
showing that time-reversal invariance is
equivalent to
\begin{equation}
\boxed{\Theta H\Theta^{-1}=H.}
\end{equation}
The relation is sometimes written
$[\Theta,H]=0$, but $\Theta$
is anti\emph{linear}; the equation
above is the precise statement.
\textbf{(c)} If $H\psi_n=E_n\psi_n$
with real $E_n$, then
$H\Theta\psi_n=\Theta H\psi_n
=E_n\Theta\psi_n$, so
$\psi_n^*$ has the same energy.
For a nondegenerate eigenvalue,
$\psi_n^*=c\psi_n$, $|c|=1$.
Choose a global phase
$\psi_n\mapsto\ee^{\ii\alpha}\psi_n$
to make $c=1$, proving that an
eigenfunction can be chosen real.
\textbf{(d)} Momentum eigenstates satisfy
\[
\boxed{\Theta f_p(x)=f_p^*(x)=f_{-p}(x)}.
\]
Both energies equal $p^2/(2m)$,
although their momenta are opposite.
For hydrogen
$\psi_{n\ell m}=R_{n\ell}(r)Y_\ell^m$,
using the real radial convention and
$Y_\ell^{m*}=(-1)^mY_\ell^{-m}$ gives
\begin{equation}
\boxed{\Theta\psi_{n\ell m}
=(-1)^m\psi_{n\ell,-m}.}
\end{equation}
Hydrogen energies depend only
on $n$ in the nonrelativistic
point-Coulomb model, so the
time-reversed state has the same
energy. For $m=0$ it is the same
state up to phase, while for
$m\ne0$ it is generally an
independent, degenerate state.
\phys{A fixed externally imposed
magnetic field breaks this symmetry
unless the magnetic field itself
is reversed with the system.}

\problem{6.37}{Time reversal of spin-half and Kramers degeneracy}
A spin-$1/2$ spinor is
$\chi=(a,b)^T$. With the
Condon--Shortley spin basis,
time reversal acts as
\begin{equation}
\boxed{\Theta\begin{pmatrix}a\\b\end{pmatrix}
=\begin{pmatrix}-b^*\\a^*\end{pmatrix}
=-\ii\sigma_yK\begin{pmatrix}a\\b\end{pmatrix}.}
\end{equation}
\textbf{(a)} Apply the operation a second time:
\[
\Theta^2\begin{pmatrix}a\\b\end{pmatrix}
=\Theta\begin{pmatrix}-b^*\\a^*\end{pmatrix}
=\begin{pmatrix}-a\\-b\end{pmatrix}.
\]
Therefore
\begin{equation}
\boxed{\Theta^2=-I\qquad(s=\tfrac12).}
\end{equation}
\textbf{(b)} Suppose a normalized energy
eigenstate $\ket\psi$ of a
time-reversal invariant Hamiltonian
were identical, up to phase,
to its time-reversed state:
$\Theta\ket\psi=c\ket\psi$.
As $\Theta$ is anti\emph{linear},
applying it again gives
\[
\Theta^2\ket\psi
=\Theta(c\ket\psi)
=c^*\Theta\ket\psi
=|c|^2\ket\psi.
\]
But part (a) demands
$\Theta^2\ket\psi=-\ket\psi$.
Thus $|c|^2=-1$, impossible.
Hence $\ket\psi$ and
$\Theta\ket\psi$ are distinct
eigenstates with the same energy:
\begin{equation}
\boxed{\text{Every level is at least doubly
degenerate if }\Theta^2=-I
\text{ and }[\Theta,H]=0.}
\end{equation}
The two states are in fact orthogonal.
Using antiunitarity
$\braket{\Theta a}{\Theta b}=
\braket b a$ with $a=\psi$, $b=\Theta\psi$,
\[
\braket{\Theta\psi}{\Theta^2\psi}
=\braket{\Theta\psi}{\psi}.
\]
The left side equals
$-\braket{\Theta\psi}{\psi}$,
so their overlap is zero.
This is Kramers degeneracy.
\caveat{The degeneracy theorem
applies to systems whose physical
time-reversal operator squares to
$-I$, such as an odd number of
spin-half fermions under the stated
assumptions. A generic spinless
particle has $\Theta^2=+I$ and
need not have a Kramers doublet.}
